\documentclass[11pt]{article}
\usepackage[margin=2.0cm]{geometry}
\usepackage{booktabs}
\usepackage{amsmath}
\usepackage{amssymb}
\usepackage{enumitem}
\usepackage{array}
\usepackage{listings}
\usepackage{xcolor}
\usepackage[T1]{fontenc}
\usepackage{lmodern}

\setlength{\parindent}{0pt}
\setlength{\parskip}{3.6pt}
\newcommand{\code}[1]{\texttt{#1}}
\newcommand{\hr}{\vspace{2pt}\hrule\vspace{6pt}}
\renewcommand{\arraystretch}{1.12}

\definecolor{kw}{RGB}{0,90,160}
\definecolor{cmt}{RGB}{110,110,110}
\lstdefinestyle{vhdl}{
  language=VHDL, basicstyle=\ttfamily\footnotesize,
  keywordstyle=\color{kw}\bfseries, commentstyle=\color{cmt}\itshape,
  showstringspaces=false, columns=fullflexible, breaklines=true,
  frame=leftline, framesep=6pt, xleftmargin=8pt, aboveskip=6pt, belowskip=6pt
}
\lstset{style=vhdl}

\begin{document}

\begin{center}
{\Large\bfseries EE\,362 --- Computer Organization Laboratory}\\[2pt]
{\large\bfseries Lab 6 --- Guided Examples: Solution}\\[2pt]
{\itshape Ex 1 a file small enough to watch \quad$\cdot$\quad Ex 2 the second read port \quad$\cdot$\quad Ex 3 read-during-write and \$0}\\[4pt]
Fall 2026 \quad$\cdot$\quad Lab TA: Gandhi Torres
\end{center}
\hr

\begin{center}
\small
\begin{tabular}{cl>{\raggedright\arraybackslash}p{6.2cm}>{\raggedright\arraybackslash}p{3.6cm}}
\toprule
& circuit & new idea & feeds the assignment as \\
\midrule
1 & \code{regfile\_small} ($4 \times 8$) & storage as an \textbf{array}, clocked write, combinational read & the shape of \code{RegFile} \\
2 & \code{regfile\_two\_ports} & \textbf{read ports are cheap} & \code{Aout} and \code{Bout} \\
3 & \code{read\_during\_write} & two decisions that look like details & read-old, and \$0 \\
\bottomrule
\end{tabular}
\end{center}

\vspace{4pt}\hr

\section*{Example 1 --- A register file small enough to watch}

Four registers, eight bits, one read port. The real thing is $32 \times 32$ with
two read ports: \textbf{the same design}, and impossible to see in a waveform.

\subsection*{1.1 Three components you already own}

\begin{center}
\begin{tabular}{lll}
\toprule
part & from & question it answers \\
\midrule
address \textbf{decoder} & Lab 0, Ex 3 & which register does this address mean? \\
array of \textbf{registers} & Lab 2 & where is the value stored? \\
\textbf{multiplexer} & Lab 3 & whose output reaches the port? \\
\bottomrule
\end{tabular}
\end{center}

\begin{quote}
\textbf{A register file is not a new kind of component.} That is the whole
lesson of Example 1, and it is why this lab is buildable in one session.
\end{quote}

\subsection*{1.2 The array type}

\begin{lstlisting}
type reg_array is array (0 to 3) of std_logic_vector(7 downto 0);
signal regs : reg_array;
\end{lstlisting}

\code{regs(2)} is the third register. \textbf{The index must be an integer}, which
is why the address arrives as bits and gets converted --- Lab 4, Example 3:

\begin{lstlisting}
regs(to_integer(unsigned(Dsel))) <= Din;
\end{lstlisting}

\subsection*{1.3 The asymmetry: write is clocked, reads are not}

\begin{center}
\begin{tabular}{lll}
\toprule
& when it happens & why \\
\midrule
\textbf{write} & on the rising edge & storage needs an edge to be unambiguous \\
\textbf{reads} & continuously & change the address and the output follows at once \\
\bottomrule
\end{tabular}
\end{center}

\begin{quote}
\textbf{That asymmetry is deliberate, and it is exactly what the CPU needs.} In
Lab 12 the register file must produce both operands \emph{during} a cycle and
store the result \emph{at the end} of it.
\end{quote}

So the reads sit \textbf{outside} the clocked process. Putting them inside is the
most common structural mistake in this lab: it delays every read by a cycle and
the CPU then executes each instruction with last cycle's operands.

\vspace{4pt}\hr

\section*{Example 2 --- The second read port, and why it is nearly free}

Compare Example 2 with Example 1. \textbf{The difference is two lines:}

\begin{lstlisting}
Bsel : in std_logic_vector(1 downto 0);          -- one more address port
Bout <= regs(to_integer(unsigned(Bsel)));        -- one more multiplexer
\end{lstlisting}

\textbf{The storage did not change at all.}

\subsection*{2.1 Read ports versus write ports}

\begin{center}
\small
\begin{tabular}{l>{\raggedright\arraybackslash}p{11.0cm}}
\toprule
another \textbf{read} port & one more multiplexer hanging off storage that already exists. Any number of readers can look at a register at once, so there is no conflict to resolve. \\
\midrule
another \textbf{write} port & two sources may target the \textbf{same register on the same edge}. The hardware must decide who wins: arbitration, priority, and a mux on every register's input. \textbf{It is a new correctness problem, not one more mux.} \\
\bottomrule
\end{tabular}
\end{center}

Compare the resource numbers of Examples 1 and 2 in your Fitter report to see the
read-port cost yourself. It is small.

\subsection*{2.2 Why exactly two reads and one write}

Look at any R-type instruction: \code{ADD rd, rs, rt}. Two source registers are
read, one destination is written.

\begin{quote}
\textbf{The register file's shape is dictated by the instruction format}, not
chosen for convenience. That is also why a second write port is not needed: no
instruction in this machine writes two registers.
\end{quote}

\vspace{4pt}\hr

\section*{Example 3 --- Two decisions that look like details}

\subsection*{3.1 Read-during-write}

What should a read port produce when it addresses the register being written on
this very edge?

\begin{center}
\begin{tabular}{lll}
\toprule
& the read port returns & cost \\
\midrule
\textbf{read-old} & the value stored before the edge & free --- it is what you get naturally \\
\textbf{read-new} (bypass) & the value being written now & a comparator and a mux \emph{per read port} \\
\bottomrule
\end{tabular}
\end{center}

\textbf{There is no universally right answer.} Both are buildable and both are
used in real machines.

\begin{quote}
\textbf{Your design is read-old unless you deliberately add the bypass}, because
the read multiplexer looks at each register's current output and nothing is stored
until the edge.
\end{quote}

That is correct for this CPU: in Lab 12 the register file is read at the start of
a cycle and written at the end, so the two never collide in a way that matters.

\textbf{What you must do:} know which one you built, and say so in the report. A
student who cannot answer ``old or new?'' has not understood their own design ---
and it is a standard interview question.

\textbf{How to test it:} put a known value in register $k$; then on one edge write
a \emph{different} value to $k$ while reading $k$. Sample just before the edge and
just after. The two designs disagree only on that one sample.

\subsection*{3.2 Register zero}

In MIPS, \code{\$0} always reads as zero --- not by a convention software politely
follows, but because \textbf{the hardware refuses to store into it}:

\[
\mathit{We}_k = \code{Strb} \cdot m_k(\code{Dsel}) \cdot \underbrace{[\,k \neq 0\,]}_{\text{the whole mechanism}}
\]

\begin{lstlisting}
if Strb = '1' and Dsel /= "00000" then
    regs(to_integer(unsigned(Dsel))) <= Din;
end if;
\end{lstlisting}

\begin{quote}
\textbf{``Hardwired to zero'' is not a rule. It is a missing wire.}
\end{quote}

\textbf{Protect the read side as well:}

\begin{lstlisting}
Aout <= (others => '0') when Asel = "00000"
   else regs(to_integer(unsigned(Asel)));
\end{lstlisting}

A design that blocks only the write passes its first test --- the array starts at
zeros --- and returns stale data the moment anything targets register 0.

\textbf{Why have such a register at all?} It makes the instruction set smaller: a
permanent zero gives you several instructions for free. A move is an add with
\code{\$0}; a clear is a write of \code{\$0}; a comparison against zero needs no
literal.

\vspace{4pt}\hr

\section*{Into the assignment}

\begin{center}
\begin{tabular}{ll}
\toprule
example & becomes \\
\midrule
Ex 1 & the array, the conversion, the clocked write and combinational reads \\
Ex 2 & two read ports, \code{Aout} and \code{Bout}, independent addresses \\
Ex 3 & read-old, and register 0 blocked on \textbf{both} sides \\
\bottomrule
\end{tabular}
\end{center}

Four things to check in your own \code{RegFile}:

\begin{enumerate}[leftmargin=2em,itemsep=2pt]
\item register 0 blocked on the \textbf{write} side (\code{Dsel /= "00000"});
\item register 0 forced on the \textbf{read} side too;
\item reads \textbf{outside} the clocked process;
\item \textbf{no \code{else}} on the \code{Strb} test --- not writing means hold,
the same rule as Labs 1 and 2.
\end{enumerate}

\textbf{Expect this to be the largest component so far.} Thirty-two 32-bit
registers is 1024 flip-flops, and in Lab 12's resource table the register file
dominates the whole CPU. That is worth one sentence in your report, because it is
the first time a design decision shows up clearly as area.

\vspace{6pt}
\hr
{\small \textbf{Forward.} In Lab 12 this component supplies both ALU operands and
receives every result. Its two read addresses come from the instruction's
\code{rs} and \code{rt} fields, and its write address from \code{rd} or
\code{rt} --- chosen by a multiplexer you built in Lab 3.}

\end{document}
